\documentclass[11pt]{article}
\usepackage[a4paper,margin=1.9cm]{geometry}
\usepackage[T1]{fontenc}
\usepackage{textcomp}
\usepackage[spanish]{babel}
\usepackage{amsmath,amssymb}
\usepackage{lmodern}
\renewcommand{\familydefault}{\sfdefault}
\usepackage{enumitem}

\newcommand{\vect}[2]{\begin{pmatrix}#1\\#2\end{pmatrix}}
\usepackage{tikz}
\usetikzlibrary{calc,arrows.meta,patterns,decorations.pathmorphing}
\usepackage[table]{xcolor}
\usepackage{multicol}
\usepackage{parskip}

\definecolor{unalblue}{RGB}{31,73,124}

\newlist{opciones}{enumerate}{1}
\setlist[opciones]{label=(\alph*),itemsep=6pt,topsep=6pt,leftmargin=1.8em,font=\normalfont}
\setlist[enumerate,1]{itemsep=1.4em,topsep=1em}

\newcommand{\examfooter}{%
  \vfill
  \rule{\linewidth}{0.4pt}\\[1pt]
  {\scriptsize\textit{Transcripción digital --- MathUNAL}\hfill\textcolor{unalblue}{\textbf{\thepage}}}
}
\pagestyle{empty}

\newcommand{\nota}[1]{{\small\itshape\color{unalblue!70!black} #1}}

\begin{document}
\noindent
\begin{minipage}[t]{0.6\linewidth}
{\small\bfseries Geometría Vectorial y Analítica --- Parcial 1}\\
{\small Semestre 2016--01}
\end{minipage}%
\hfill
\begin{minipage}[t]{0.35\linewidth}
\raggedleft\small Escuela de Matemáticas. Universidad Nacional de Colombia, Sede Medellín.
\end{minipage}\\[2pt]
\rule{\linewidth}{0.6pt}

\noindent
\begin{minipage}[t]{0.54\linewidth}
\textbf{Primer Examen Parcial de Geometría}\\
12 de marzo del 2016\\[6pt]
\textbf{Puntaje. Sólo para uso oficial}

\small
\begin{tabular}{|c|c|c|c|c|c|}
\hline
1 (/30) & 2 (/30) & 3 (/10) & 4-6 (/30) & Total & Nota \\
\hline
 & & & & & \\
\hline
\end{tabular}
\end{minipage}%
\begin{minipage}[t]{0.44\linewidth}
Nombre: \underline{\hspace{4.8cm}}\\[8pt]
Cédula: \underline{\hspace{1.8cm}}\ \ Salón: \underline{\hspace{1.2cm}}\\[8pt]
Grupo: \underline{\hspace{1.8cm}}\ \ Examen No.: \underline{\hspace{1.6cm}}\\[8pt]
Firma: \underline{\hspace{4.8cm}}
\end{minipage}

\vspace{10pt}
\textbf{Instrucciones:} La duración del examen es de 1 hora y 40 minutos. El examen consta de ocho preguntas en tres hojas impresas por ambos lados, antes de empezar verifique que su examen esté \textbf{completo} y que sus \textbf{datos} sean \textbf{correctos}. No \textbf{desengrape} su exámen ni \textbf{añada} otras grapas o su examen será \textbf{anulado}. En las preguntas con procedimiento \textbf{justifique} sus respuestas en los \textbf{espacios asignados}. No está permitido hacer preguntas, el uso de calculadoras, sacar hojas en blanco ni ningún tipo de apuntes durante el examen. Por favor verifique que su celular esté \textbf{apagado}.

\vspace{6pt}
\nota{En cada uno de los literales de los puntos 1 y 2, debe presentar detalladamente el procedimiento para llegar a la solución. NOTA: Respuestas sin procedimiento no se tendrán en cuenta.}

\begin{enumerate}
\item Considere la recta $\mathcal{L}$ dada por la ecuación general $x-2y+10=0$.

\begin{enumerate}[label=(\alph*)]
\item{}[10pt] Encuentre unas ecuaciones escalares paramétricas para la recta $\mathcal{L}$.
\item{}[10pt] Considere la recta $\mathcal{L}'$ que es perpendicular a $\mathcal{L}$ y pasa por el punto $P=\vect{1}{1}$. Encuentre el punto de intersección entre $\mathcal{L}$ y $\mathcal{L}'$.
\item{}[10pt] Encuentre el punto $Q$ de la recta $\mathcal{L}$, que está más cercano al origen.
\end{enumerate}

\item Considere la recta $\mathcal{L}$ dada por la ecuación general $\sqrt{3}x-y-\sqrt{3}=0$.

\begin{enumerate}[label=(\alph*)]
\item{}[10pt] Encuentre el ángulo de inclinación de $\mathcal{L}$.
\item{}[10pt] Encuentre la ecuación de la bisectriz del ángulo agudo formado por $\mathcal{L}$ y el eje $x$. \textbf{Observación}: Recuerde que la bisectriz de un ángulo es la recta que divide al ángulo en dos ángulos iguales.
\item{}[10pt] Halle dos vectores unitarios, normales a la recta $\mathcal{L}$.
\end{enumerate}

\item{}[10pt] Sean $X$ y $U$ dos vectores no nulos y suponga que $\|X+U\|=\|X-U\|$. Demuestre que $X$ y $U$ son ortogonales.
\end{enumerate}

\vspace{6pt}
\nota{En las preguntas 4 a 6 complete los espacios en blanco o elija la (las) opción (opciones) correcta (correctas) según el caso. Tenga presente que los datos en cada literal son independientes y pueden variar. NOTA: En esta sección se califica sólo la respuesta y no se tiene en cuenta el procedimiento.}

\begin{enumerate}
\setcounter{enumi}{3}
\item{}[10pt] En la figura se muestra la recta $\mathcal{L}$, con pendiente $m$, que pasa por el origen, así como también los vectores $U$, $V$ y $W$. Responda a las siguientes preguntas.

\begin{center}
\begin{tikzpicture}[scale=0.9]
\draw[-{Latex}] (-3,0) -- (3,0) node[right]{$x$};
\draw[-{Latex}] (0,-2.5) -- (0,2.5) node[left]{$y$};
\draw[thick] (-2.5,-1.5) -- (2.5,1.5) node[below right]{$\mathcal{L}$};
\draw[-{Latex},thick] (0,0) -- (1,2) node[right]{$U=\vect{1}{2}$};
\draw[-{Latex},thick] (0,0) -- (-2,-2) node[below left]{$V=\vect{v_1}{-2}$};
\draw[-{Latex},thick] (0,0) -- (-2.2,1.6) node[above left]{$W=\vect{-1}{w_2}$};
\end{tikzpicture}
\end{center}

\begin{enumerate}[label=(\alph*)]
\item{}[3pt] Si $\|V\|=2\sqrt{2}$ entonces $v_1=$ \underline{\hspace{4cm}}.
\item{}[3pt] Si $U$ es normal a $\mathcal{L}$ entonces $m=$ \underline{\hspace{4cm}}.
\item{}[4pt] Si $U\cdot W=5$ entonces $\|W\|=$ \underline{\hspace{4cm}}.
\end{enumerate}

\item[5.] [10pt] Las rectas $\mathcal{L}_1$, $\mathcal{L}_2$, $\mathcal{L}_3$ de la figura tienen ángulos de inclinación $\alpha_1,\alpha_2,\alpha_3$ y pendientes $m_1,m_2,m_3$ respectivamente. Se entiende que $\mathcal{L}_1$ y $\mathcal{L}_2$ son perpendiculares.

\begin{center}
\begin{tikzpicture}[scale=0.9]
\draw[-{Latex}] (-3,0) -- (3.5,0) node[right]{$x$};
\draw[-{Latex}] (0,-1.5) -- (0,2.5) node[above]{$y$};
\draw[thick] (-2.2,-1.6) -- (2,1.5) node[right]{$\mathcal{L}_2$};
\draw[thick] (-1,-1.6) -- (1.6,1.9) node[above]{};
\draw[thick] (-3,1.2) -- (3,-0.9);
\node at (2.6,-0.75) {$\mathcal{L}_1$};
\node at (0.9,-0.15) {$\mathcal{L}_3$};
\node[font=\small] at (-1.4,-0.15) {$\alpha_1$};
\node[font=\small] at (0.3,0.35) {$\theta$};
\node[font=\small] at (0.55,-0.15) {$\alpha_3$};
\node[font=\small] at (1.5,-0.15) {$\alpha_2$};
\end{tikzpicture}
\end{center}

\begin{enumerate}[label=(\alph*)]
\item{}[3pt] Si $m_2=-2$ entonces $\tan\alpha_1=$ \underline{\hspace{4cm}}.
\item{}[4pt] Si $\mathcal{L}_1: -2x+3y-2=0$ y además $\mathcal{L}_3=\left\{t\vect{1}{3}+\vect{2}{0}: t\in\mathbb{R}\right\}$, entonces $\cos\theta=$ \underline{\hspace{4cm}}.
\item{}[3pt] Si $\tan\alpha_2=-1$ y $m_3=3$ entonces $\tan\theta=$ \underline{\hspace{4cm}}.
\end{enumerate}

\nota{Indicación. Se recuerda que $\tan(\beta+\gamma)=\dfrac{\tan\beta+\tan\gamma}{1-\tan\beta\tan\gamma}$ y $\tan(\beta-\gamma)=\dfrac{\tan\beta-\tan\gamma}{1+\tan\beta\tan\gamma}$.}

\item[6.] [10pt] En la figura a continuación, las rectas $\mathcal{L}_1$ y $\mathcal{L}_3$ son ortogonales, el vector $N$ es su punto de intersección y dirige a $\mathcal{L}_3$. La recta $\mathcal{L}_2$ pasa por el origen y es paralela a $\mathcal{L}_1$, $U$ es un vector del plano.

\begin{center}
\begin{tikzpicture}[scale=0.9]
\draw[-{Latex}] (-3,0) -- (3,0) node[right]{$x$};
\draw[-{Latex}] (0,-1.5) -- (0,2.5) node[left]{$y$};
\draw[thick] (-2.5,-1.2) -- (2,2.4) node[above]{$\mathcal{L}_3$};
\draw[thick] (-2.7,-1.9) -- (2.5,1.5) node[right]{$\mathcal{L}_1$};
\draw[thick] (-2.7,1.35) -- (2.7,-1.35) node[below right]{$\mathcal{L}_2$};
\draw[-{Latex},thick] (0,0) -- (0.9,1.9) node[right]{$U=\vect{u_1}{u_2}$};
\node[left,font=\small] at (-0.3,0.15) {$N=\vect{3}{2}$};
\end{tikzpicture}
\end{center}

\begin{enumerate}[label=(\alph*)]
\item{}[3pt] Si $\mathrm{P}_{\mathcal{L}_3}U=\vect{6}{4}$ entonces $\operatorname{dist}(U,\mathcal{L}_1)=$ \underline{\hspace{4cm}}.
\item{}[3pt] Si $\|\mathrm{P}_{\mathcal{L}_2}U\|=3$ y $\|U\|=4$ entonces $\operatorname{dist}(U,\mathcal{L}_2)=$ \underline{\hspace{4cm}}.
\item{}[4pt] Suponga que $Y=t\vect{-3}{-2}+U$ con $t\in\mathbb{R}$ y además $\operatorname{dist}(Y,\mathcal{L}_3)=2$ entonces $\operatorname{dist}(U,\mathcal{L}_3)=$ \underline{\hspace{4cm}}.
\end{enumerate}

\end{enumerate}

\examfooter
\end{document}
